\documentclass{situationist-preprint}

\title{Constitutive Rules of Speech Acts}
\msid{LNG-13-a}
\msfamily{LNG-13}
\msdiscipline{Philosophy of Language}
\mstype{Principal manuscript}
\mskeywords{speech acts; illocution; constitutive rules; Austin; Searle; Strawson; assertion; promising}

\begin{document}
\maketitle

\begin{abstract}
What is a speech act, and what makes two speech acts acts of the same kind? Austin distinguished the act performed in saying something from the act of saying it and from its effects. Searle held that illocutionary acts are constituted by rules, as moves in a game are; Strawson and Bach and Harnish held that most of them are individuated by the communicative intentions with which they are performed. We argue that a speech act is an act performed in uttering something, and that it is individuated by its constitutive rules. Perlocutionary effects cannot individuate illocutionary acts, since the act is complete when the effect fails; communicative intentions cannot, since the same intentions are present in acts of different kinds and where no act is performed. The individuating rules are not the conventions of a particular language, and they fix what an act is, not whether it binds.
\end{abstract}

\section{The problem}

\citet{Austin1962} distinguished three things a speaker does. In uttering \enquote{The bull is about to charge}, she says something with a sense and reference, a locutionary act; in saying it, she may warn her hearer, an illocutionary act; by saying it, she may alarm him, a perlocutionary act. Austin held that illocutionary acts are conventional. \citet{Searle1969} developed this into the thesis that illocutionary acts are constituted by rules of the form \enquote{$X$ counts as $Y$ in context $C$}, as a move in chess is constituted by the rules of chess.

\citet{Strawson1964} objected that the conventional element is essential only to a restricted class of illocutions, those tied to institutions, such as baptizing, bidding or pronouncing a verdict. Most communicative illocutions, such as warning and asserting, are individuated by Gricean intentions that are meant to be recognized \citep{Grice1957}. \citet{BachHarnish1979} developed the view systematically.

We argue for Searle's side, in a form that meets Strawson's objection.

\section{Framework}

\begin{definition}[Speech act]\label{d:sa}
A \emph{speech act} is an act performed \emph{in} uttering something: an illocutionary act in Austin's sense, as distinct from the locutionary act and from perlocutionary acts performed \emph{by} uttering.
\end{definition}

\begin{definition}[Constitutive and regulative rules]\label{d:rules}
A \emph{regulative rule} governs a form of activity that exists independently of it (\enquote{drive on the right}). A \emph{constitutive rule} creates the possibility of a form of activity by specifying that something counts as something else in a context (\enquote{moving the king out of check counts as a legal reply}) \citep{Rawls1955,Searle1969}.
\end{definition}

\begin{definition}[Individuation]\label{d:ind}
A property $\Phi$ \emph{individuates} speech acts if two speech acts are of the same type exactly when they share $\Phi$, and an utterance performs an act of a type exactly when it has the $\Phi$ of that type.
\end{definition}

\begin{assumption}[Correctness]\label{a:correct}
Speech acts of different types have different conditions of correct performance (an assertion and a guess that $p$ are criticizable in different circumstances), and an act of a type is criticizable as such when it violates its conditions.
\end{assumption}

\section{The argument}

\begin{lemma}\label{l:perl}
Perlocutionary effects do not individuate speech acts.
\end{lemma}

\begin{proof}
A warning that does not alarm its hearer is still a warning, and an argument that does not persuade is still an argument; and a hearer can be alarmed by a warning, a threat or a report. By Definition~\ref{d:ind}, the effect does not individuate.
\end{proof}

\begin{lemma}\label{l:int}
Communicative intentions do not individuate speech acts.
\end{lemma}

\begin{proof}
First, the same communicative intentions can be present in acts of different types. A speaker who guesses that $p$ and a speaker who asserts that $p$ may each intend her hearer to take her to believe $p$, and to recognize that intention; yet by Assumption~\ref{a:correct} the assertion is criticizable if she does not know that $p$, and the guess is not. Second, the same intentions can be present where no act of the type is performed. A bystander who says \enquote{I name this ship the \emph{Queen Elizabeth}}, with every intention that the appointed person has, has not named the ship; the appointed person, saying the same words, has \citep{Austin1962}. By Definition~\ref{d:ind}, communicative intentions do not individuate.
\end{proof}

\begin{theorem}\label{t:main}
A speech act is an act performed in uttering something, individuated by its constitutive rules.
\end{theorem}

\begin{proof}
The first conjunct is Definition~\ref{d:sa}. For the second, by Assumption~\ref{a:correct}, speech-act types differ in their conditions of correct performance, and an utterance is criticizable as an act of a type when it violates the conditions of that type. An utterance can violate the conditions of a type only if it is an act of that type, so its being of that type cannot depend on its satisfying them. What makes it an act of that type is that it counts as such under a rule specifying what counts as performing it and what conditions the performance answers to, and such a rule is constitutive (Definition~\ref{d:rules}). Such rules individuate in the sense of Definition~\ref{d:ind}: two utterances that count as the same $Y$ under the rules perform the same act, and an utterance that counts as a $Y$ performs a $Y$. By Lemmas~\ref{l:perl} and~\ref{l:int}, the two alternatives that might individuate speech acts instead do not.
\end{proof}

The knowledge rule for assertion proposed by \citet{Williamson2000} is a rule of this kind: an assertion that breaks it is still an assertion, and faulty as one.

\begin{corollary}\label{c:promise}
The constitutive rules of promising determine what counts as a promise; they do not by themselves make a promise morally binding.
\end{corollary}

\begin{proof}
By Theorem~\ref{t:main}, the rules of promising make an utterance of the right form, in the right circumstances, count as an undertaking. That the speaker has undertaken something is a fact about how her utterance counts under the rules; that she is morally obligated to do it does not follow from such facts without a normative premise, as MTE-10 shows for \enquote{ought} and \enquote{is} generally.
\end{proof}

Corollary~\ref{c:promise} is where \citet{Searle1964}'s derivation of \enquote{ought} from \enquote{is} fails, and it is what ETH-16 requires: constitutive rules make an utterance a promise, while what makes a promise bind is the promisor's creation of reasonable reliance in the promisee.

\section{Objections and replies}

\begin{objection}[Non-conventional illocutions]\label{o:strawson}
A warning can be given by any words, or by none, in a language that has no convention for warning. If illocutionary acts were constituted by rules, they would require conventions, and most do not \citep{Strawson1964}.
\end{objection}

\begin{reply}
Constitutive rules are not conventions. \citet{Searle1969} distinguished the rules of an act type from the conventional devices by which different languages realize them: the rule that an utterance counts as an undertaking is common to English and French, while the words used to perform it differ. An act is individuated by its rule, and any device, conventional or not, by which a speaker makes her utterance count under the rule suffices to perform it. This fits ETH-16, on which a social practice may supply the means of promising but is not the source of a promise's force.
\end{reply}

\begin{objection}[Expressed attitudes]\label{o:bh}
The guess and the assertion differ in the attitude expressed, a hunch in one case and a belief in the other, so intentions to express attitudes do individuate them \citep{BachHarnish1979}.
\end{objection}

\begin{reply}
To express a belief rather than a hunch is not to have a different intention about the hearer's uptake; it is to make one's utterance answerable to a different condition. That is a status conferred by a rule, so the reply concedes the thesis.
\end{reply}

\begin{objection}[Unknown rules]\label{o:regress}
Speakers cannot state the rules of assertion or promising. Rules they cannot state cannot be the rules they follow, so cannot individuate their acts.
\end{objection}

\begin{reply}
A rule can be followed without being formulated, and to require that it be stated as a premise of the speaker's reasoning would start the regress that LOG-18 shows to arise when rules are treated as premises.
\end{reply}

\section{Conclusion}

A speech act is an act performed in uttering something. It is individuated neither by its effects nor by communicative intentions. It is individuated by the constitutive rules under which the utterance counts as an act of its kind and is answerable to that kind's conditions.

\begin{thebibliography}{99}\small
\bibitem[Austin(1962)]{Austin1962} Austin, J.~L. (1962). \emph{How to Do Things with Words}. Clarendon Press.
\bibitem[Bach and Harnish(1979)]{BachHarnish1979} Bach, K. and Harnish, R.~M. (1979). \emph{Linguistic Communication and Speech Acts}. MIT Press.
\bibitem[Grice(1957)]{Grice1957} Grice, H.~P. (1957). Meaning. \emph{The Philosophical Review}, 66(3), 377--388.
\bibitem[Rawls(1955)]{Rawls1955} Rawls, J. (1955). Two concepts of rules. \emph{The Philosophical Review}, 64(1), 3--32.
\bibitem[Searle(1964)]{Searle1964} Searle, J.~R. (1964). How to derive \enquote{ought} from \enquote{is}. \emph{The Philosophical Review}, 73(1), 43--58.
\bibitem[Searle(1969)]{Searle1969} Searle, J.~R. (1969). \emph{Speech Acts: An Essay in the Philosophy of Language}. Cambridge University Press.
\bibitem[Strawson(1964)]{Strawson1964} Strawson, P.~F. (1964). Intention and convention in speech acts. \emph{The Philosophical Review}, 73(4), 439--460.
\bibitem[Williamson(2000)]{Williamson2000} Williamson, T. (2000). \emph{Knowledge and Its Limits}. Oxford University Press.
\end{thebibliography}

\end{document}
